\section{The Electromagnetic Field Tensor}
\label{sec:electromagnetic_field_tensor}

The electric and magnetic fields are not separate relativistic objects. A Lorentz transformation mixes them, just as a boost mixes the scalar and vector potentials. Their unified description is the antisymmetric electromagnetic field tensor.

\subsection{Definition}

The electromagnetic field tensor is defined by

\[
\boxed{
F_{\mu\nu}
=
\partial_{\mu}A_{\nu}
-
\partial_{\nu}A_{\mu}
}.
\]

Raising both indices gives

\[
\boxed{
F^{\mu\nu}
=
\partial^{\mu}A^{\nu}
-
\partial^{\nu}A^{\mu}
}.
\]

The two formulas are equivalent because the Minkowski metric is constant in Cartesian inertial coordinates.

\subsection{Antisymmetry}

Interchanging the indices gives

\[
\begin{aligned}
F_{\nu\mu}
&=
\partial_{\nu}A_{\mu}
-
\partial_{\mu}A_{\nu}\\
&=
-F_{\mu\nu}.
\end{aligned}
\]

Therefore

\[
\boxed{
F_{\mu\nu}
=
-F_{\nu\mu}
}.
\]

In particular,

\[
F_{\mu\mu}=0
\]

with no sum over \(\mu\).

A general four-by-four matrix has sixteen components. Antisymmetry reduces the number of independent components to

\[
\frac{4(4-1)}{2}=6.
\]

These six components correspond exactly to the three components of \(\mathbf{E}\) and the three components of \(\mathbf{B}\).

\subsection{Gauge invariance}

Under

\[
A'_{\mu}
=
A_{\mu}
-
\partial_{\mu}\chi,
\]

the transformed field tensor is

\[
\begin{aligned}
F'_{\mu\nu}
&=
\partial_{\mu}A'_{\nu}
-
\partial_{\nu}A'_{\mu}\\
&=
\partial_{\mu}
\left(
A_{\nu}-\partial_{\nu}\chi
\right)
-
\partial_{\nu}
\left(
A_{\mu}-\partial_{\mu}\chi
\right)\\
&=
F_{\mu\nu}
-
\partial_{\mu}\partial_{\nu}\chi
+
\partial_{\nu}\partial_{\mu}\chi.
\end{aligned}
\]

For a smooth gauge function, partial derivatives commute. Hence

\[
\boxed{
F'_{\mu\nu}=F_{\mu\nu}
}.
\]

The field tensor is gauge invariant even though the four-potential is not.

\subsection{Lorentz transformation law}

Because \(A^{\mu}\) is a four-vector, the field tensor transforms as a rank-two tensor:

\[
\boxed{
F'^{\mu\nu}
=
\Lambda^{\mu}{}_{\rho}
\Lambda^{\nu}{}_{\sigma}
F^{\rho\sigma}
}.
\]

This transformation mixes components identified as electric in one frame with components identified as magnetic in another frame.

The tensor is therefore the natural relativistic electromagnetic field variable.

\subsection{Exterior-derivative interpretation}

Introduce the potential one-form

\[
A
=
A_{\mu}\,\mathrm{d}x^{\mu}.
\]

Its exterior derivative is

\[
F
=
\mathrm{d}A.
\]

In components,

\[
F
=
\frac{1}{2}
F_{\mu\nu}
\,\mathrm{d}x^{\mu}
\wedge
\mathrm{d}x^{\nu}.
\]

The antisymmetry of \(F_{\mu\nu}\) is built into the wedge product.

This geometric notation also makes the homogeneous Maxwell equations immediate:

\[
\mathrm{d}F
=
\mathrm{d}^{2}A
=0.
\]

The component form of this identity will be developed in Section~\ref{sec:homogeneous_maxwell_equations}.

\subsection{Cyclic derivative identity}

Starting from the definition,

\[
\begin{aligned}
&\partial_{\lambda}F_{\mu\nu}
+
\partial_{\mu}F_{\nu\lambda}
+
\partial_{\nu}F_{\lambda\mu}
\\
&=
\partial_{\lambda}
\left(
\partial_{\mu}A_{\nu}
-
\partial_{\nu}A_{\mu}
\right)
+
\partial_{\mu}
\left(
\partial_{\nu}A_{\lambda}
-
\partial_{\lambda}A_{\nu}
\right)
\\
&\quad+
\partial_{\nu}
\left(
\partial_{\lambda}A_{\mu}
-
\partial_{\mu}A_{\lambda}
\right).
\end{aligned}
\]

Every second-derivative term cancels with another term after commuting partial derivatives. Thus

\[
\boxed{
\partial_{\lambda}F_{\mu\nu}
+
\partial_{\mu}F_{\nu\lambda}
+
\partial_{\nu}F_{\lambda\mu}
=0
}.
\]

This is the tensor form of the Bianchi identity for electromagnetism.

\subsection{Dual field tensor}

Define the Levi--Civita symbol by

\[
\epsilon^{0123}=+1.
\]

The dual electromagnetic tensor is

\[
\boxed{
\widetilde{F}^{\mu\nu}
=
\frac{1}{2}
\epsilon^{\mu\nu\rho\sigma}
F_{\rho\sigma}
}.
\]

It is also antisymmetric:

\[
\widetilde{F}^{\mu\nu}
=
-\widetilde{F}^{\nu\mu}.
\]

The cyclic identity is equivalent to

\[
\boxed{
\partial_{\mu}
\widetilde{F}^{\mu\nu}
=0
}.
\]

This compact equation contains both homogeneous Maxwell equations.

\subsection{Independent tensor components}

Because of antisymmetry, one may choose the six independent components as

\[
F_{01},
\quad
F_{02},
\quad
F_{03},
\quad
F_{12},
\quad
F_{23},
\quad
F_{31}.
\]

The time-space components encode the electric field, while the purely spatial components encode the magnetic field.

The precise signs depend on

\begin{itemize}
    \item the metric signature;
    \item the definition of the four-potential;
    \item whether indices on \(F\) are raised or lowered.
\end{itemize}

The next section derives the component matrices explicitly for the conventions used in these notes.

\subsection{Lorentz scalar formed from the tensor}

The contraction

\[
F_{\mu\nu}F^{\mu\nu}
\]

is a Lorentz scalar. Because \(F_{\mu\nu}\) is gauge invariant, this contraction is also gauge invariant.

The free electromagnetic Lagrangian density in SI units will be

\[
\boxed{
\mathcal{L}_{\mathrm{EM}}
=
-
\frac{1}{4\mu_0}
F_{\mu\nu}F^{\mu\nu}
}.
\]

The factor \(1/4\) compensates for double counting caused by antisymmetry. The variation of this action will produce the inhomogeneous Maxwell equations in the next chapter.

\subsection{Why an antisymmetric derivative?}

A simple derivative \(\partial_{\mu}A_{\nu}\) is not gauge invariant. Under a gauge transformation it acquires a second derivative of \(\chi\).

The antisymmetric combination removes that term:

\[
\partial_{\mu}A_{\nu}
-
\partial_{\nu}A_{\mu}.
\]

This construction is the simplest local first-derivative expression that is both Lorentz covariant and gauge invariant.

\subsection{Reconstruction of the potential}

If a tensor \(F_{\mu\nu}\) satisfies the cyclic identity in a contractible region, then locally there exists a potential such that

\[
F_{\mu\nu}
=
\partial_{\mu}A_{\nu}
-
\partial_{\nu}A_{\mu}.
\]

The potential is determined only up to a gauge transformation.

The local statement can fail globally on regions with nontrivial topology. This is the tensor version of the topological qualification encountered for scalar and vector potentials.

\subsection{Common mistakes}

Typical errors include:

\begin{enumerate}
    \item defining \(F_{\mu\nu}\) without antisymmetrizing the derivatives;
    \item forgetting that raising a spatial index can change a sign;
    \item treating the potential as gauge invariant because the field tensor is gauge invariant;
    \item counting sixteen independent components instead of six;
    \item using the Levi--Civita symbol without stating the convention for \(\epsilon^{0123}\);
    \item assuming that the existence of a local potential guarantees a single global potential on every domain.
\end{enumerate}

\subsection{What has been established}

The electromagnetic field is represented covariantly by

\[
F_{\mu\nu}
=
\partial_{\mu}A_{\nu}
-
\partial_{\nu}A_{\mu}.
\]

The tensor is antisymmetric, Lorentz covariant, and gauge invariant. It satisfies the cyclic identity because it is constructed from a potential. The next section identifies its six independent components with the electric and magnetic fields.
